\addcontentsline{toc}{section}{References}
\begin{thebibliography}{99}
\small
\setlength{\itemsep}{3pt}
\setlength{\parskip}{0pt}

\bibitem{ciocan} B. I. Ciocan et al. (2026),
``MUSE-DARK III: The evolution of the radial acceleration relation at
intermediate redshifts,'' \emph{Astronomy \& Astrophysics} 709, L16.
\href{https://arxiv.org/abs/2604.22613}{arXiv:2604.22613}.
The coefficient comparison uses Appendix E's self-consistent inference.

\bibitem{unruh} W. G. Unruh (1976), ``Notes on black-hole evaporation,''
\emph{Physical Review D} 14, 870.
\href{https://doi.org/10.1103/PhysRevD.14.870}{doi:10.1103/PhysRevD.14.870}.

\bibitem{gibbons} G. W. Gibbons and S. W. Hawking (1977),
``Cosmological event horizons, thermodynamics, and particle creation,''
\emph{Physical Review D} 15, 2738.
\href{https://doi.org/10.1103/PhysRevD.15.2738}{doi:10.1103/PhysRevD.15.2738}.

\bibitem{milgrom1999} M. Milgrom (1999), ``The modified dynamics as a vacuum effect,''
\emph{Physics Letters A} 253, 273.
\href{https://doi.org/10.1016/S0375-9601(99)00077-8}{doi:10.1016/S0375-9601(99)00077-8}.

\bibitem{warren} H. T. Warren et al. (2025), ``Estimate of Current Mass
of the Large Magellanic Cloud from the Orphan-Chenab Tidal Stream.''
\href{https://arxiv.org/abs/2508.01086}{arXiv:2508.01086}.

\bibitem{brooks} R. A. N. Brooks et al. (2026), ``Quantifying the
Milky Way, LMC and their interaction using all-sky kinematics of outer
halo stars.'' \href{https://arxiv.org/abs/2603.08788}{arXiv:2603.08788}.

\bibitem{lelli} F. Lelli et al. (2015), ``Gas dynamics in tidal dwarf
galaxies: disc formation at $z=0$.''
\href{https://arxiv.org/abs/1509.05404}{arXiv:1509.05404}.

\bibitem{wolf} J. Wolf et al. (2010), ``Accurate masses for
dispersion-supported galaxies,'' \emph{MNRAS} 406, 1220.
\href{https://arxiv.org/abs/0908.2995}{arXiv:0908.2995}.

\bibitem{read19} J. I. Read, M. G. Walker, and P. Steger (2019), ``Dark
matter heats up in dwarf galaxies,'' \emph{MNRAS} 484, 1401.

\bibitem{moster} B. P. Moster, T. Naab, and S. D. M. White (2013),
``Galactic star formation and accretion histories from matching galaxies
to dark matter haloes.''
\href{https://arxiv.org/abs/1205.5807}{arXiv:1205.5807}.

\bibitem{dutton} A. A. Dutton and A. V. Macci\`o (2014),
``Cold dark matter haloes in the Planck era: evolution of structural
parameters for Einasto and NFW profiles.''
\href{https://arxiv.org/abs/1402.7073}{arXiv:1402.7073}.

\bibitem{atzberger} K. R. Atzberger et al. (2026),
``The DECam MAGIC Survey: Uncovering the Tidal Tails of the Crater II
Dwarf Galaxy.''
\href{https://arxiv.org/abs/2602.21283}{arXiv:2602.21283}.

\bibitem{limberg} G. Limberg et al. (2025),
``$S^5$: Tidal Disruption in Crater 2 and Formation of Diffuse Dwarf
Galaxies in the Local Group.''
\href{https://arxiv.org/abs/2512.02177}{arXiv:2512.02177}.

\bibitem{ji} A. P. Ji et al. (2021),
``Kinematics of Antlia 2 and Crater 2 from The Southern Stellar Stream
Spectroscopic Survey ($S^5$).''
\href{https://arxiv.org/abs/2106.12656}{arXiv:2106.12656}.

\bibitem{desi} DESI Collaboration (2025),
``DESI DR2 Results II: Measurements of Baryon Acoustic Oscillations and
Cosmological Constraints.''
\href{https://arxiv.org/abs/2503.14738}{arXiv:2503.14738}.

\bibitem{katz96} N. Katz, D. H. Weinberg and L. Hernquist (1996),
``Cosmological simulations with TreeSPH,''
\emph{Astrophys. J. Suppl.} \textbf{105}, 19--35.

\bibitem{thoul96} A. A. Thoul and D. H. Weinberg (1996),
``Hydrodynamic simulations of galaxy formation. II. Photoionization and
the formation of low-mass galaxies,'' \emph{Astrophys. J.} \textbf{465}, 608.

\bibitem{okamoto08} T. Okamoto, L. Gao and T. Theuns (2008),
``Mass loss of galaxies due to an ultraviolet background,''
\emph{Mon. Not. R. Astron. Soc.} \textbf{390}, 920--928.

\bibitem{oshea15} B. W. O'Shea, J. H. Wise, H. Xu and M. L. Norman (2015),
``Probing the ultraviolet luminosity function of the earliest galaxies with the Renaissance Simulations,''
\emph{Astrophys. J. Lett.} \textbf{807}, L12.

\bibitem{xu16} H. Xu, J. H. Wise, M. L. Norman, K. Ahn and B. W. O'Shea (2016),
``Galaxy properties and UV escape fractions during the epoch of reionization: results from the Renaissance Simulations,''
\emph{Astrophys. J.} \textbf{833}, 84.

\bibitem{nadler20} E. O. Nadler et al. (2020),
``Milky Way satellite census. II. Galaxy--halo connection constraints including the impact of the Large Magellanic Cloud,''
\emph{Astrophys. J.} \textbf{893}, 48.

\bibitem{planck} Planck Collaboration (2020),
``Planck 2018 results. VI. Cosmological parameters.''
\href{https://arxiv.org/abs/1807.06209}{arXiv:1807.06209}.

\bibitem{desrecal} B. Popovic et al. (2025),
``The Dark Energy Survey Supernova Program: A Reanalysis Of Cosmology
Results And Evidence For Evolving Dark Energy With An Updated Type Ia
Supernova Calibration.''
\href{https://arxiv.org/abs/2511.07517}{arXiv:2511.07517}.

\bibitem{desrepo} DES Science, public supernova data and covariance
repository. \href{https://github.com/des-science/DES-SN5YR}{DES-SN5YR on
GitHub}. The archived data version used by the calculations should be
preserved when reproducing them.

\bibitem{releasecode} J. Chinitz, public dark-energy reproduction release.
\href{https://github.com/jakechinitz/dark-energy/releases/tag/Dark-Energy}
{Dark-Energy release on GitHub}. The current mixing extension is
documented by local sources C1--C3.


\bibitem{bcgrar} Brightest-cluster-galaxy radial acceleration relation,
\href{https://arxiv.org/abs/2402.12016}{arXiv:2402.12016}.

\bibitem{centralsrar} Radial acceleration relation of central galaxies in
groups and clusters, \href{https://arxiv.org/abs/2603.23591}{arXiv:2603.23591}.

\bibitem{angus08} G. W. Angus, B. Famaey and D. A. Buote (2008), ``X-ray
group and cluster mass profiles in MOND,'' \emph{MNRAS} 387, 1470.
\href{https://arxiv.org/abs/0709.0108}{arXiv:0709.0108}.

\bibitem{brouwer21} M. M. Brouwer et al. (2021), ``The weak lensing radial
acceleration relation,'' \emph{A\&A} 650, A113.
\href{https://arxiv.org/abs/2106.11677}{arXiv:2106.11677}.

\bibitem{lelli17} F. Lelli, S. S. McGaugh, J. M. Schombert and M. S. Pawlowski
(2017), ``One law to rule them all: the radial acceleration relation of
galaxies,'' \emph{ApJ} 836, 152.
\href{https://arxiv.org/abs/1610.08981}{arXiv:1610.08981}.

\bibitem{humphrey11} P. J. Humphrey et al. (2011), ``A census of baryons and
dark matter in an isolated, Milky Way sized elliptical galaxy,'' \emph{ApJ}
729, 53. \href{https://arxiv.org/abs/1010.6078}{arXiv:1010.6078}.

\bibitem{julio25} M. P. J\'ulio et al. (2025), ``The radial acceleration
relation at the EDGE of galaxy formation: testing its universality in
low-mass dwarf galaxies,'' \emph{A\&A} 704, A330.

\bibitem{mayer22} A. C. Mayer, A. F. Teklu, K. Dolag and R.-S. Remus (2022),
``$\Lambda$CDM with baryons vs.\ MOND: the time evolution of the universal
acceleration scale in the Magneticum simulations,''
\href{https://arxiv.org/abs/2206.04333}{arXiv:2206.04333}.

\bibitem{xcopnew} S. Ettori et al., ``Hydrostatic mass profiles in X-COP galaxy clusters,'' A\&A 621, A39 (2019).
\href{https://arxiv.org/abs/1805.00035}{arXiv:1805.00035};
\href{https://dominiqueeckert.wixsite.com/xcop/data}{public X-COP data}.
\bibitem{herbonnetnew} R. Herbonnet et al., ``CCCP and MENeaCS: (updated) weak-lensing masses for 100 galaxy clusters,'' MNRAS 497, 4684 (2020), Table A2.
\href{https://arxiv.org/abs/1912.04414}{arXiv:1912.04414}.
\bibitem{bulletnew} G. Rihtar\v{s}i\v{c} et al., ``Mapping dark matter in the Bullet Cluster using JWST imaging and spectroscopy'' (2026).
\href{https://arxiv.org/html/2601.22245v2}{arXiv:2601.22245v2};
\href{https://niriss.github.io/lensing.html}{CANUCS Bullet v1 maps}.
\bibitem{sarkarnew} A. Sarkar et al., ``XRISM Observations of Abell 1795: Evidence for Low Turbulence and Resonant Scattering'' (2026).
\href{https://arxiv.org/abs/2606.08097}{arXiv:2606.08097}.
\bibitem{xrismnew} XRISM Collaboration, ``Constraining gas motion and non-thermal pressure beyond the core of the Abell 2029 galaxy cluster with XRISM'' (2025).
\href{https://arxiv.org/abs/2505.06533}{arXiv:2505.06533}.


\bibitem{sparcnew} SPARC collaboration, \href{https://astroweb.case.edu/SPARC/}{public rotation-curve and baryonic-mass data}. Selections and nuisance priors are recorded in U1, S03.
\bibitem{kidsnew} KiDS collaboration, \href{https://kids.strw.leidenuniv.nl/sciencedata.php}{public science data}; T. Mistele et al., \href{https://arxiv.org/abs/2310.15248}{arXiv:2310.15248}, updated acceleration table. Covariance conventions are recorded in U1, S04.
\bibitem{actnew} ACT collaboration, \href{https://github.com/ACTCollaboration/DR6-ACT-lite}{DR6-ACT-lite} and \href{https://lambda.gsfc.nasa.gov/product/act/act_dr6.02/index.html}{DR6 public products}; lensing analysis \href{https://arxiv.org/abs/2304.05202}{arXiv:2304.05202} and \href{https://arxiv.org/abs/2304.05203}{arXiv:2304.05203}.
\bibitem{egnew} ACT--BOSS lensing--velocity analysis, \href{https://arxiv.org/abs/2405.12795}{arXiv:2405.12795}. Band and redshift conventions are recorded in U1, S09.
\bibitem{constantsnew} NIST, \href{https://physics.nist.gov/cuu/Constants/}{CODATA reference values}; Particle Data Group, \href{https://pdglive.lbl.gov/DataBlock.action?node=S035M}{tau mass data}. The archived comparison values are recorded in U1, S11.
\bibitem{gwtcnew} LIGO--Virgo--KAGRA, \href{https://dcc.ligo.org/LIGO-P2600128/public}{public GWTC-4 gravity-test products}. The polarization evidence combination is recorded in U1, S12.

\bibitem{gastaldello07} F. Gastaldello et al. (2007), ``Probing the dark
matter and gas fraction in relaxed galaxy groups with X-ray observations from
Chandra and XMM-Newton,'' \emph{ApJ} 669, 158.
\href{https://arxiv.org/abs/astro-ph/0610134}{arXiv:astro-ph/0610134}.
\bibitem{vikhlinin06} A. Vikhlinin et al. (2006), ``Chandra sample of nearby
relaxed galaxy clusters: mass, gas fraction, and mass-temperature relation,''
\emph{ApJ} 640, 691. \href{https://arxiv.org/abs/astro-ph/0507092}{arXiv:astro-ph/0507092}.
\end{thebibliography}
